\documentclass[11pt]{article}
\usepackage[margin=1in]{geometry}
\usepackage{amsmath,amssymb,amsthm,mathtools}
\usepackage{booktabs,tabularx,array,longtable}
\usepackage{xcolor}
\usepackage{enumitem}
\usepackage{hyperref}
\usepackage{microtype}
\usepackage{tikz-cd}
\usepackage{fancyhdr}
\usepackage{titlesec}
\usepackage{newtxtext,newtxmath}

\definecolor{navy}{RGB}{35,55,85}
\definecolor{softblue}{RGB}{235,242,250}
\definecolor{softgray}{RGB}{246,246,246}
\definecolor{softgreen}{RGB}{237,247,240}
\definecolor{softgold}{RGB}{251,246,232}
\definecolor{softred}{RGB}{251,238,238}
\hypersetup{colorlinks=true,linkcolor=navy,citecolor=navy,urlcolor=navy}
\titleformat{\section}{\Large\bfseries\color{navy}}{\thesection}{0.7em}{}
\titleformat{\subsection}{\large\bfseries\color{navy}}{\thesubsection}{0.7em}{}
\pagestyle{fancy}
\setlength{\headheight}{14pt}
\fancyhf{}
\lhead{ProLT Observation Refinement}
\rhead{Classification note v0.2}
\cfoot{\thepage}
\setlength{\parindent}{0pt}
\setlength{\parskip}{0.55em}

\newtheorem{theorem}{Theorem}[section]
\newtheorem{proposition}[theorem]{Proposition}
\newtheorem{lemma}[theorem]{Lemma}
\newtheorem{corollary}[theorem]{Corollary}
\theoremstyle{definition}
\newtheorem{definition}[theorem]{Definition}
\newtheorem{example}[theorem]{Example}
\newtheorem{remark}[theorem]{Remark}

\newcommand{\F}{\mathbb F_2}
\newcommand{\Om}{\Omega_P}
\newcommand{\truth}{\varsigma}
\newcommand{\OX}{\mathcal O}
\newcommand{\Cone}{\operatorname{Cone}}
\newcommand{\im}{\operatorname{im}}
\newcommand{\rank}{\operatorname{rank}}
\newcommand{\xor}{\mathbin{\Updownarrow}}
\newcommand{\redhom}{\widetilde H}
\newcommand{\statusS}{\textbf{[S]}}
\newcommand{\statusA}{\textbf{[A]}}
\newcommand{\statusC}{\textbf{[C]}}
\newcommand{\down}{\mathord\downarrow}
\newcommand{\up}{\mathord\uparrow}

\begin{document}

\begin{center}
{\LARGE\bfseries\color{navy} Observation Refinement in Propositional Logical Topologies}\par
\vspace{0.3em}
{\large Fiber profiles, adjoint criteria, Quillen fibers, refinement defects, and the limits of complete classification}\par
\vspace{0.65em}
{\normalsize Working research note v0.2 --- 23 September 2026}\par
\vspace{0.3em}
{\small Follow-up to \emph{Homological Foundations for Propositional Logical Topologies}, built on the ProLT core manuscript v0.8.}
\end{center}

\vspace{0.7em}
\noindent\colorbox{softgold}{\parbox{0.96\textwidth}{
\textbf{Status convention.} \statusS\ denotes standard finite-poset/algebraic-topological mathematics. \statusA\ denotes a direct ProLT specialization or elementary theorem derived from the v0.8 signature framework. \statusC\ denotes a potentially publishable ProLT-specific formulation, invariant, or algorithmic consequence whose novelty should be subjected to a dedicated literature audit before any priority claim.
}}

\section{Question and answer in one page}

Let \(\Phi=(\phi_i)_{i\in I}\) be an observation family and let \(\Psi\) extend \(\Phi\) by new observations. The v0.8 ProLT theory gives the realized-signature posets
\[
P_\Phi=o_\Phi(\Om),\qquad P_\Psi=o_\Psi(\Om),
\]
ordered by inclusion, and a canonical forgetful refinement map
\[
r_{\Psi\Phi}:P_\Psi\longrightarrow P_\Phi
\]
obtained by deleting the new coordinates. The corresponding order complexes are
\[
\OX_\Phi=\Delta(P_\Phi),\qquad \OX_\Psi=\Delta(P_\Psi).
\]
The central problem is:

\begin{center}
\emph{When does adding observations change only the resolution of the presentation, and when does it change the homotopy or homology of the observable-state space?}
\end{center}

The main results of this note give a hierarchy of answers.

\begin{enumerate}[label=\textbf{R\arabic*.},leftmargin=2.7em]
\item Every refinement admits an exact \emph{fiber-profile normal form}: over each coarse signature \(p\), record exactly which new-coordinate signatures occur. This completely reconstructs \(P_\Psi\) and the map \(r_{\Psi\Phi}\).
\item The refinement is \emph{topologically redundant} exactly when every coarse state has one new profile and those profiles are monotone. Equivalently, \(r_{\Psi\Phi}\) is an order isomorphism.
\item If the refinement does not split coarse states, then \(\OX_\Psi\) is literally a subcomplex of \(\OX_\Phi\). A coarse chain survives exactly when its new-coordinate profiles increase along the chain. For one new Boolean observation, the deleted chains are precisely those containing a \(1\to0\) descent.
\item There are simple signature-level sufficient criteria for homotopy invariance. If every lower refinement fiber has a greatest element, or every upper refinement fiber has a least element, then \(r_{\Psi\Phi}\) has an order adjoint and is a homotopy equivalence. For one observation \(\psi\), this says: \emph{possibility of \(\psi\)} or \emph{forced truth of \(\psi\)} must be monotone in the coarse information order.
\item More generally, Quillen--McCord/Barmak fiber criteria apply. Contractible principal fibers imply a simple homotopy equivalence; acyclic principal fibers imply homology equivalence.
\item Homology equivalence has an exact finite test: the refinement mapping cone is acyclic. In the non-splitting case it reduces to ordinary relative homology \(H_*(\OX_\Phi,\OX_\Psi)\).
\item There is no unrestricted algorithmic classification of homotopy-neutral refinements: arbitrary finite posets can be realized as refined signature posets above a one-point coarse space, so the problem contains finite-complex contractibility, which is undecidable. Homology neutrality, by contrast, remains decidable by finite boundary-matrix computation.
\end{enumerate}

\section{Base ProLT data and refinement map}

Let \(P=\{p_1,\dots,p_n\}\) be a finite propositional variable set and \(\Om=\{0,1\}^P\) its valuation space. For a formula \(L\), write
\[
\truth(L)=\{v\in\Om:v\models L\}.
\]
For an indexed observation family \(\Phi=(\phi_i)_{i\in I}\), define
\[
o_\Phi(v)=\{i\in I:v\models\phi_i\}\subseteq I.
\]
The v0.8 ProLT results identify topological indistinguishability with equality of signatures and the specialization preorder with inclusion:
\[
v\sim_\Phi w\iff o_\Phi(v)=o_\Phi(w),
\qquad
v\preceq_\Phi w\iff o_\Phi(v)\subseteq o_\Phi(w).
\]
Hence the Kolmogorov quotient is the finite poset
\[
P_\Phi=o_\Phi(\Om)\subseteq 2^I
\]
with inclusion order.

Suppose now that
\[
\Psi=(\Phi,\Theta)
\]
extends \(\Phi\) by a disjointly indexed family \(\Theta=(\theta_a)_{a\in A}\). Then
\[
o_\Psi(v)=\bigl(o_\Phi(v),o_\Theta(v)\bigr)
\in 2^I\times 2^A,
\]
and the coordinate-forgetting map
\[
r=r_{\Psi\Phi}:P_\Psi\to P_\Phi,
\qquad
r(s)=s\cap I,
\]
is surjective and monotone.

\section{The fiber-profile normal form}

\begin{definition}[Coarse state and refinement profile]\label{def:profile}
For \(p\in P_\Phi\), let
\[
C_p:=o_\Phi^{-1}(p)
\]
be the coarse observational state represented by \(p\). Define its \emph{new-coordinate fiber profile}
\[
\mathcal B_p:=o_\Theta(C_p)
=\{o_\Theta(v):v\in C_p\}\subseteq 2^A.
\]
Each \(\mathcal B_p\) is nonempty.
\end{definition}

\begin{theorem}[Fiber-profile normal form]\label{thm:normalform}
\statusA/\statusC\;
Let
\[
Q_{\mathcal B}
:=\{(p,a):p\in P_\Phi,\ a\in\mathcal B_p\}
\subseteq P_\Phi\times2^A
\]
with product order
\[
(p,a)\le(q,b)
\iff p\le q\text{ in }P_\Phi\text{ and }a\subseteq b.
\]
Then
\[
P_\Psi\cong Q_{\mathcal B}
\]
as posets, under separation of old and new coordinates. Under this isomorphism the refinement map is exactly the first-coordinate projection
\[
r(p,a)=p.
\]
Consequently the family \((\mathcal B_p)_{p\in P_\Phi}\) determines the refined signature poset and the refinement projection completely.
\end{theorem}

\begin{proof}
Every realized \(\Psi\)-signature comes from a valuation \(v\) and splits uniquely as
\[
(o_\Phi(v),o_\Theta(v))=(p,a)
\]
with \(a\in\mathcal B_p\). Conversely, every pair in \(Q_{\mathcal B}\) is realized by the definition of \(\mathcal B_p\). Inclusion of the combined signatures is coordinatewise inclusion because the old and new index sets are disjoint. The projection statement is immediate.
\end{proof}

\begin{remark}[Point fibers are not Quillen fibers]\label{rem:pointfiber}
The direct point fiber
\[
r^{-1}(p)=\{p\}\times\mathcal B_p
\]
records only how the coarse state \(p\) splits. Quillen-type theorems use instead the preimages of principal ideals or filters,
\[
r^{-1}(\down p),\qquad r^{-1}(\up p),
\]
which also remember how the split pieces compare with states below or above \(p\). This distinction is essential: every point fiber can be contractible while the refinement still changes homotopy type.
\end{remark}

\section{Exact classification of topological redundancy}

The v0.8 manuscript calls a new observation topologically redundant when its truth region is already open in the old ProLT. The fiber normal form turns this into an exact signature-level theorem for an arbitrary finite block of added observations.

\begin{theorem}[Redundancy/order-isomorphism theorem]\label{thm:redundancy}
\statusA/\statusC\;
For a refinement \(\Psi=(\Phi,\Theta)\), the following are equivalent.
\begin{enumerate}[label=(\roman*)]
\item \(\tau_\Psi=\tau_\Phi\) on the fixed valuation carrier \(\Om\).
\item Every new truth region \(\truth(\theta_a)\), \(a\in A\), is already \(\tau_\Phi\)-open.
\item Every fiber profile is a singleton,
\[
\mathcal B_p=\{\alpha(p)\},
\]
and the resulting map
\[
\alpha:P_\Phi\to2^A
\]
is isotone:
\[
p\le q\Longrightarrow\alpha(p)\subseteq\alpha(q).
\]
\item The refinement projection \(r:P_\Psi\to P_\Phi\) is an order isomorphism.
\end{enumerate}
\end{theorem}

\begin{proof}
The equivalence of (i) and (ii) is the definition of equality of the generated topologies: \(\tau_\Phi\subseteq\tau_\Psi\) automatically, and equality holds exactly when each new subbasic truth region is old-open.

Assume (i). If two valuations lie in the same old signature class \(C_p\), they are topologically indistinguishable in \(\tau_\Phi\); equality of topologies makes them indistinguishable in \(\tau_\Psi\) as well. Thus the new signature is constant on \(C_p\), giving \(\mathcal B_p=\{\alpha(p)\}\). If \(p\le q\), choose \(v\in C_p\), \(w\in C_q\). The old specialization relation gives \(v\preceq_\Phi w\). Equality of topologies gives \(v\preceq_\Psi w\), hence \(\alpha(p)\subseteq\alpha(q)\). Therefore (iii) holds.

Assume (iii). For each new coordinate \(a\in A\), the set
\[
U_a:=\{p\in P_\Phi:a\in\alpha(p)\}
\]
is upward closed. By the signature representation of the ProLT topology,
\[
\truth(\theta_a)=o_\Phi^{-1}(U_a)
\]
is \(\tau_\Phi\)-open. Hence (ii).

Under (iii), Theorem~\ref{thm:normalform} identifies \(P_\Psi\) with the graph \(\{(p,\alpha(p))\}\). Monotonicity of \(\alpha\) says
\[
p\le q\iff(p,\alpha(p))\le(q,\alpha(q)),
\]
so projection is an order isomorphism. Conversely, if projection is an order isomorphism then every fiber has one point and its inverse is necessarily isotone, giving (iii).
\end{proof}

\begin{remark}
This theorem separates three notions which should not be conflated:
\begin{enumerate}[label=(\alph*)]
\item \emph{no splitting}: each \(\mathcal B_p\) is a singleton;
\item \emph{no order deletion}: the singleton profile map is isotone;
\item \emph{full redundancy}: both hold, equivalently the topology is unchanged.
\end{enumerate}
A refinement can leave every coarse equivalence class intact while still deleting one-way specialization comparisons.
\end{remark}

\section{Non-splitting refinements and chain deletion}

\begin{definition}[Non-splitting refinement]
A refinement is \emph{non-splitting} if every \(\mathcal B_p\) is a singleton. Write
\[
\mathcal B_p=\{\alpha(p)\}.
\]
No monotonicity of \(\alpha\) is assumed.
\end{definition}

\begin{theorem}[Order deletion theorem]\label{thm:orderdelete}
\statusA/\statusC\;
For a non-splitting refinement, identify the fine and coarse vertex sets through the bijection \(r\). Then
\[
p\le_\Psi q
\iff
p\le_\Phi q\ \text{and}\ \alpha(p)\subseteq\alpha(q).
\]
Consequently
\[
\le_\Psi\ \subseteq\ \le_\Phi
\]
and
\[
\OX_\Psi\subseteq\OX_\Phi
\]
as a simplicial subcomplex on the same vertex set. A coarse chain
\[
p_0<p_1<\cdots<p_k
\]
survives the refinement exactly when
\[
\alpha(p_0)\subseteq\alpha(p_1)\subseteq\cdots\subseteq\alpha(p_k).
\]
\end{theorem}

\begin{proof}
In the fiber-profile normal form the unique fine point over \(p\) is \((p,\alpha(p))\). Product order gives the displayed criterion. Every fine chain is therefore a coarse chain; conversely, a coarse chain is fine exactly when its new profiles are coordinatewise nondecreasing.
\end{proof}

\begin{corollary}[One-bit descent criterion]\label{cor:descent}
Suppose one new Boolean observation \(\psi\) is constant on each coarse class, and write
\[
b(p)\in\{0,1\}
\]
for its value on \(C_p\). Then a coarse chain \(p_0<\cdots<p_k\) survives exactly when
\[
b(p_0)\le b(p_1)\le\cdots\le b(p_k).
\]
Equivalently, it is deleted exactly when it contains a \(1\to0\) inversion.
\end{corollary}

\begin{corollary}[Post-\(T_0\) refinement filtration]\label{cor:T0filtration}
\statusA/\statusC\;
If \(X_\Phi\) is already \(T_0\) on the valuation carrier, then \(o_\Phi\) is injective, so every later observation refinement is automatically non-splitting. Thus a sequence
\[
\Phi_0\subseteq\Phi_1\subseteq\cdots\subseteq\Phi_T
\]
with \(X_{\Phi_0}\) \(T_0\) yields, after identifying vertices with valuations,
\[
\OX_{\Phi_T}\subseteq\cdots\subseteq\OX_{\Phi_1}\subseteq\OX_{\Phi_0}.
\]
Hence ordinary persistence can be applied after reversing the refinement index; zigzag persistence is not required while observations are only added.
\end{corollary}

\begin{proof}
A coarse signature class contains a single valuation when \(o_\Phi\) is injective, so it cannot split further as a set. Theorem~\ref{thm:orderdelete} then applies at every later stage.
\end{proof}

\section{Relative defect chains in the non-splitting case}

The preceding inclusion lets the mapping-cone defect from the foundation note be sharpened.

\begin{theorem}[Refinement defect becomes relative homology]\label{thm:relativecone}
\statusS/\statusA\;
For a non-splitting refinement, under the vertex identification of Theorem~\ref{thm:orderdelete}, the induced simplicial map is the inclusion
\[
\OX_\Psi\hookrightarrow\OX_\Phi.
\]
Therefore the mapping-cone refinement defect satisfies
\[
D_k(\Psi/\Phi;R)
\cong
H_k(\OX_\Phi,\OX_\Psi;R).
\]
\end{theorem}

\begin{definition}[Refinement-descent chain complex]\label{def:descentcomplex}
For a non-splitting refinement define \(C_k^{\mathrm{desc}}(\Psi/\Phi;R)\) to be the free \(R\)-module generated by the coarse \(k\)-chains that fail the survival condition
\[
\alpha(p_0)\subseteq\cdots\subseteq\alpha(p_k).
\]
Its boundary is the ordinary simplicial boundary followed by deletion of every face that does survive the refinement.
\end{definition}

\begin{proposition}[Exact descent-chain computation]\label{prop:descentcomplex}
\statusA/\statusC\;
There is a natural chain isomorphism
\[
C_*^{\mathrm{desc}}(\Psi/\Phi;R)
\cong
C_*(\OX_\Phi,\OX_\Psi;R).
\]
Hence
\[
D_k(\Psi/\Phi;R)
\cong
H_k(C_*^{\mathrm{desc}}(\Psi/\Phi;R)).
\]
For one added Boolean observation, the generators are exactly coarse chains containing a \(1\to0\) descent.
\end{proposition}

\begin{proof}
The simplicial chain group of \(\OX_\Phi\) is free on its simplices, and the subcomplex \(\OX_\Psi\) is spanned by exactly the surviving simplices. The quotient chain group therefore has the nonsurviving simplices as a basis. Its differential is the coarse simplicial differential modulo the surviving submodule, which is precisely the stated rule.
\end{proof}

\section{Adjoint criteria: an elementary homotopy-neutral class}

For \(p\in P_\Phi\), define the principal refinement fibers
\[
L_p:=r^{-1}(\down p),
\qquad
U_p:=r^{-1}(\up p).
\]
The following test is much easier than computing arbitrary homotopy types.

\begin{theorem}[Greatest-lower-fiber / least-upper-fiber criterion]\label{thm:adjoint}
\statusS/\statusA/\statusC\;
Let
\[
\mathsf U(p)
:=\bigcup_{\substack{q\le p\\a\in\mathcal B_q}}a,
\qquad
\mathsf L(p)
:=\bigcap_{\substack{q\ge p\\a\in\mathcal B_q}}a.
\]
Then:
\begin{enumerate}[label=(\roman*)]
\item \(L_p\) has a greatest element if and only if
\[
\mathsf U(p)\in\mathcal B_p.
\]
When this holds, the greatest element is \((p,\mathsf U(p))\).
\item \(U_p\) has a least element if and only if
\[
\mathsf L(p)\in\mathcal B_p.
\]
When this holds, the least element is \((p,\mathsf L(p))\).
\end{enumerate}
If (i) holds for every \(p\), then \(r\) has a right adjoint. If (ii) holds for every \(p\), then \(r\) has a left adjoint. Either global condition implies that
\[
\Delta(r):\OX_\Psi\to\OX_\Phi
\]
is a homotopy equivalence.
\end{theorem}

\begin{proof}
Suppose \(L_p\) has greatest element \((q,a_*)\). Since \(r\) is surjective, choose some \((p,b)\in r^{-1}(p)\subseteq L_p\). Greatness gives \((p,b)\le(q,a_*)\), so \(p\le q\le p\) and therefore \(q=p\). The second coordinate \(a_*\) contains every new signature occurring over every \(q\le p\), hence contains \(\mathsf U(p)\). Since \(a_*\) itself occurs in the union, equality follows. Conversely, if \(\mathsf U(p)\in\mathcal B_p\), then every \((q,a)\in L_p\) satisfies
\[
q\le p,\qquad a\subseteq\mathsf U(p),
\]
so \((p,\mathsf U(p))\) is greatest. This proves (i). The proof of (ii) is dual.

If every lower fiber has a greatest element, define
\[
s_R(p):=(p,\mathsf U(p)).
\]
Then
\[
r(x)\le p\iff x\le s_R(p),
\]
so \(r\dashv s_R\). In particular \(r s_R=\mathrm{id}\) and
\[
\mathrm{id}\le s_Rr.
\]
Order-preserving maps that are pointwise comparable induce homotopic maps on order complexes, so \(\Delta(s_R)\) is a homotopy inverse to \(\Delta(r)\). The left-adjoint case is dual.
\end{proof}

\subsection{A sharp one-observation interpretation}

Let \(\psi\) be one newly added observation. For each coarse state \(C_p\), define
\[
M_\psi(p)=
\begin{cases}
1,&\text{some }v\in C_p\text{ satisfies }\psi,\\
0,&\text{otherwise},
\end{cases}
\]
and
\[
m_\psi(p)=
\begin{cases}
1,&\text{every }v\in C_p\text{ satisfies }\psi,\\
0,&\text{otherwise}.
\end{cases}
\]
Thus \(M_\psi\) records \emph{possibility} of the new observation inside a coarse state, while \(m_\psi\) records \emph{forced truth}. The possible fiber profiles are
\[
\{0\},\qquad\{1\},\qquad\{0,1\},
\]
with
\[
m_\psi(p)=\min\mathcal B_p,
\qquad
M_\psi(p)=\max\mathcal B_p.
\]

\begin{corollary}[Possible/forced monotonicity criterion]\label{cor:possibleforced}
\statusA/\statusC\;
For a one-observation refinement:
\begin{enumerate}[label=(\roman*)]
\item \(r\) has a right adjoint if and only if \(M_\psi:P_\Phi\to\{0,1\}\) is isotone.
\item \(r\) has a left adjoint if and only if \(m_\psi:P_\Phi\to\{0,1\}\) is isotone.
\end{enumerate}
Equivalently, either of the following is sufficient for homotopy invariance:
\begin{itemize}
\item the coarse states in which \(\psi\) is \emph{possible} form an upset;
\item the coarse states in which \(\psi\) is \emph{forced} form an upset.
\end{itemize}
Topological redundancy is the stronger special case
\[
m_\psi=M_\psi=b
\]
with \(b\) isotone.
\end{corollary}

\begin{proof}
For one bit,
\[
\mathsf U(p)=\max_{q\le p}M_\psi(q).
\]
The membership condition \(\mathsf U(p)\in\mathcal B_p\) is equivalent to \(\mathsf U(p)=M_\psi(p)\), which for every \(p\) is exactly isotonicity of \(M_\psi\). Dually,
\[
\mathsf L(p)=\min_{q\ge p}m_\psi(q)
\]
and \(\mathsf L(p)\in\mathcal B_p\) for every \(p\) is exactly isotonicity of \(m_\psi\).
\end{proof}

\section{Quillen--McCord/Barmak fiber criteria}

The adjoint criterion forces a cone point in every principal fiber. Quillen's fiber theorem permits much more general contractible fibers.

\begin{theorem}[Contractible-fiber criterion]\label{thm:quillen}
\statusS/\statusA\;
If for every \(p\in P_\Phi\), the order complex
\[
\Delta\bigl(r^{-1}(\down p)\bigr)
\]
is contractible, then
\[
\Delta(r):\OX_\Psi\longrightarrow\OX_\Phi
\]
is a simple homotopy equivalence, hence a homotopy equivalence. Dually, the same conclusion follows if every
\[
\Delta\bigl(r^{-1}(\up p)\bigr)
\]
is contractible.
\end{theorem}

\begin{proof}
The lower-fiber statement is the finite-poset McCord--Quillen theorem together with Barmak's strengthening to simple homotopy equivalence. The upper-fiber statement follows by applying the theorem to opposite posets.
\end{proof}

\begin{theorem}[Acyclic-fiber homology criterion]\label{thm:acyclicfiber}
\statusS/\statusA\;
Fix coefficients \(R\). If
\[
\redhom_i\!\left(\Delta(r^{-1}(\down p));R\right)=0
\]
for every \(p\in P_\Phi\) and every \(i\ge0\), then
\[
r_*:H_i(\OX_\Psi;R)\longrightarrow H_i(\OX_\Phi;R)
\]
is an isomorphism for every \(i\). The dual upper-fiber statement also holds.
\end{theorem}

\begin{remark}[Sufficient, not necessary]\label{rem:notnecessary}
The principal-fiber conditions are powerful but not necessary. Section~\ref{sec:examples} gives a concrete ProLT refinement for which \(r\) is a homotopy equivalence even though one lower principal fiber and one upper principal fiber are each two disconnected points. Thus neither orientation of the Quillen test succeeds, despite homotopy neutrality.
\end{remark}

\section{Exact homology classification by the refinement cone}

The homotopy problem is difficult in full generality, but the homology problem has an exact algebraic answer.

\begin{definition}[Refinement mapping cone]
Let
\[
r_\#:C_*(\OX_\Psi;R)\longrightarrow C_*(\OX_\Phi;R)
\]
be the induced chain map. Define
\[
\mathsf D_*(\Psi/\Phi;R)=\Cone(r_\#)
\]
and
\[
D_k(\Psi/\Phi;R)=H_k(\mathsf D_*(\Psi/\Phi;R)).
\]
\end{definition}

\begin{theorem}[Exact homology-equivalence criterion]\label{thm:coneexact}
\statusS/\statusA/\statusC\;
The following are equivalent:
\begin{enumerate}[label=(\roman*)]
\item \(r_*:H_k(\OX_\Psi;R)\to H_k(\OX_\Phi;R)\) is an isomorphism for every \(k\);
\item
\[
D_k(\Psi/\Phi;R)=0
\qquad\text{for every }k.
\]
\end{enumerate}
Thus the refinement cone gives an exact, finite homological classification of whether the refinement is homology-neutral.
\end{theorem}

\begin{proof}
The standard mapping-cone long exact sequence is
\[
\cdots\to H_k(\OX_\Psi)
\overset{r_*}{\longrightarrow}H_k(\OX_\Phi)
\to D_k(\Psi/\Phi)
\to H_{k-1}(\OX_\Psi)\to\cdots.
\]
Exactness gives the equivalence.
\end{proof}

\begin{corollary}[Quick obstructions]
If any of the following differs between \(\OX_\Psi\) and \(\OX_\Phi\), the refinement is not homology-neutral:
\begin{enumerate}[label=(\alph*)]
\item a Betti number over the chosen field;
\item the number of connected components;
\item the Euler characteristic.
\end{enumerate}
The converses are false in general.
\end{corollary}

\subsection{An exact \(H_0\) criterion}

\begin{proposition}[Connected-component criterion]\label{prop:H0}
\statusA\;
Let \(\pi_0(P)\) denote connected components of the undirected comparability graph of a finite poset \(P\). Since \(r\) is surjective, the induced map
\[
r_*:H_0(\OX_\Psi;R)\to H_0(\OX_\Phi;R)
\]
is always surjective. It is an isomorphism if and only if, for every connected component \(C\subseteq P_\Phi\), the inverse image
\[
r^{-1}(C)
\]
is connected.
\end{proposition}

\begin{proof}
A connected component of \(P_\Psi\) maps into one component of \(P_\Phi\) because comparability paths map to comparability-or-equality paths. Surjectivity of \(r\) ensures every coarse component receives at least one fine component. On \(H_0\), each fine component generator maps to the generator of its coarse component. The map is injective exactly when no coarse component receives more than one fine component.
\end{proof}

\section{A strict hierarchy of refinement-neutrality tests}\label{sec:hierarchy}

It is useful to distinguish five levels.

\begin{center}
\begin{tabular}{cl}
\toprule
Class & Condition \\
\midrule
\(\mathsf R\) & \(r\) is an order isomorphism (topological redundancy) \\
\(\mathsf A\) & \(r\) has a left or right order adjoint \\
\(\mathsf Q\) & all lower or all upper principal refinement fibers are contractible \\
\(\mathsf H\) & \(\Delta(r)\) is a homotopy equivalence \\
\(\mathsf M\) & \(r_*\) is a homology isomorphism \\
\bottomrule
\end{tabular}
\end{center}

Standard implications and the preceding theorems give
\[
\boxed{
\mathsf R\subseteq\mathsf A\subseteq\mathsf Q\subseteq\mathsf H\subseteq\mathsf M.
}
\]
The examples below show that the first three inclusions are already strict in very small Boolean ProLTs. The universality result of Section~\ref{sec:universality} shows that \(\mathsf H\subsetneq\mathsf M\) also occurs in unrestricted ProLT refinements by importing any finite acyclic noncontractible complex.

\section{Worked Boolean examples}\label{sec:examples}

Throughout, signatures are ordered by inclusion and coefficients for displayed homology groups are \(\F\).

\subsection{Redundant: adding \(X\wedge Y\) to \((X,Y)\)}

Let
\[
\Phi=(X,Y),\qquad \psi=X\wedge Y.
\]
The coarse signatures are
\[
D=\varnothing,
\quad A=\{X\},
\quad C=\{Y\},
\quad B=\{X,Y\},
\]
forming the Boolean diamond. The new bit is
\[
b(D)=b(A)=b(C)=0,
\qquad b(B)=1.
\]
It is isotone, so Theorem~\ref{thm:redundancy} applies. The refinement is in \(\mathsf R\): the ProLT topology and canonical order complex are unchanged, although a weighted signature metric may change if the redundant coordinate is assigned positive weight.

\subsection{Adjoint-neutral but nonredundant: adding \(X\) to no observations}

Take one propositional variable \(X\), let \(\Phi\) be the empty observation family, and refine by \(\Psi=(X)\). The coarse poset is a point; the fine poset is the two-element chain
\[
0<1.
\]
The projection to a point has both a left adjoint (choose the minimum) and a right adjoint (choose the maximum), so it is a homotopy equivalence. The topology nevertheless changes from indiscrete to Sierpi\'nski. Thus this refinement lies in \(\mathsf A\setminus\mathsf R\).

\subsection{Quillen-neutral but no adjoint: adding XOR to \((X,Y)\)}

Let
\[
\Phi=(X,Y),\qquad \psi=X\xor Y.
\]
The new-bit profile on \((D,A,C,B)\) is
\[
(0,1,1,0).
\]
No coarse class splits, but the old comparisons \(A<B\) and \(C<B\) are deleted. The fine poset has \(D\) as a minimum and three elements \(A,C,B\) above it. Hence its order complex is a cone and is contractible.

The bit is not isotone, so neither the possible nor forced map is isotone; no adjoint criterion applies. However every lower Quillen fiber is contractible: over \(D,A,C\) it is a point or edge, and over \(B\) it is the entire fine poset, again a cone. Therefore the refinement lies in
\[
\mathsf Q\setminus\mathsf A.
\]

\subsection{Homotopy-neutral while both Quillen orientations fail}\label{ex:alternating}

Take
\[
\Phi=(X\wedge Y,\ X,\ X\vee Y).
\]
On the four valuations \(00,01,10,11\), the realized signatures form a four-element chain
\[
p_0<p_1<p_2<p_3
\]
with
\[
\begin{array}{c|c}
00&\varnothing\\
01&\{X\vee Y\}\\
10&\{X,X\vee Y\}\\
11&\{X\wedge Y,X,X\vee Y\}.
\end{array}
\]
Add
\[
\psi=\neg Y.
\]
Its profile along the chain is
\[
1,0,1,0.
\]
No coarse class splits. The surviving strict comparisons are
\[
p_0<p_2,
\qquad
p_1<p_2,
\qquad
p_1<p_3.
\]
Thus \(\OX_\Psi\) is a three-edge tree, while \(\OX_\Phi\) is a \(3\)-simplex. Both are nonempty and contractible, so the refinement projection is a homotopy equivalence.

Nevertheless,
\[
r^{-1}(\down p_1)=\{p_0,p_1\}
\]
is a two-point antichain, and
\[
r^{-1}(\up p_2)=\{p_2,p_3\}
\]
is also a two-point antichain. Hence the lower-fiber Quillen criterion fails and the dual upper-fiber criterion fails. This gives a concrete refinement in
\[
\mathsf H\setminus\mathsf Q.
\]

\subsection{Homology-active: adding negation to \((X)\)}

Let
\[
\Phi=(X),\qquad \Psi=(X,\neg X).
\]
The coarse signature poset is a two-element chain, so \(\OX_\Phi\) is an edge. The fine signatures \(\{X\}\) and \(\{\neg X\}\) are incomparable, so \(\OX_\Psi\) is two isolated points. Hence
\[
H_0(\OX_\Phi)\cong\F,
\qquad
H_0(\OX_\Psi)\cong\F^2.
\]
The refinement is not homology-neutral. The mapping cone gives
\[
D_1(\Psi/\Phi;\F)\cong\F,
\qquad D_0(\Psi/\Phi;\F)=0.
\]
This is the smallest example in which adding an observation splits a connected coarse information structure into disconnected fine states.

\subsection{Splitting can still be homotopy-neutral}

Let \(\Phi=(X)\) on two variables \(X,Y\), and add \(Y\). Each coarse \(X\)-state splits into the two possible \(Y\)-values, so
\[
\mathcal B_{\varnothing}=\mathcal B_{\{X\}}=\{0,1\}.
\]
The coarse poset is an edge and the fine poset is the Boolean diamond. Both are contractible. In fact both the possible map \(M_Y\equiv1\) and the forced map \(m_Y\equiv0\) are isotone, so the projection has both adjoints. This shows that state splitting by itself does not imply a topological defect.

\section{Small exhaustive verification}

To test the logical hierarchy independently of the proofs, the accompanying script exhaustively enumerates every one-bit refinement profile over every nonempty subposet of the Boolean square \(B_2\). A coarse point may have new-bit fiber \(\{0\}\), \(\{1\}\), or \(\{0,1\}\). There are \(255\) such refinement instances in total.

The verifier checks redundancy, left/right section-adjoint criteria, lower-fiber contractibility by finite beat-point reduction, lower-fiber acyclicity over \(\F\), the Betti numbers of the coarse and fine order complexes, and the mapping-cone Betti numbers. It finds exactly five logical pattern types under those tests:
\[
\begin{array}{c|r}
(\mathsf R,\mathsf A,\mathsf Q_{\downarrow},\mathsf{Acyc}_{\downarrow},\mathsf M)&\text{count}\\\hline
(0,0,0,0,0)&24\\
(0,0,0,0,1)&16\\
(0,0,1,1,1)&20\\
(0,1,1,1,1)&144\\
(1,1,1,1,1)&51
\end{array}
\]
These finite checks are not substitutes for the theorems, but they verify the example calculations and demonstrate that the sufficient conditions are genuinely non-equivalent even at very small size.

\section{Universality of refinement and the limit of complete homotopy classification}\label{sec:universality}

The preceding criteria invite the question whether one could push further to a complete effective homotopy classification. In unrestricted dimension the answer is negative.

\begin{theorem}[Universality of refined signature posets]\label{thm:universality}
\statusA/\statusC\;
Every finite poset \(Q\) is isomorphic to the realized-signature poset of some finite propositional observation family. More strongly, \(Q\) can occur as \(P_\Psi\) in a refinement
\[
\Phi\subseteq\Psi
\]
for which \(P_\Phi\) is a single point.
\end{theorem}

\begin{proof}
Take enough propositional variables that \(|\Om|\ge|Q|\), and choose a surjection
\[
h:\Om\twoheadrightarrow Q.
\]
Let the coarse observation family \(\Phi\) be empty, so \(P_\Phi\) is one point. Index new observations by the elements \(q\in Q\). For each \(q\), prescribe the truth set
\[
T_q:=\{v\in\Om:q\le h(v)\}.
\]
Every subset of a finite valuation space is definable by a propositional formula, for example by the disjunction of the complete minterms of its members. Choose \(\theta_q\) with \(\truth(\theta_q)=T_q\).

If \(h(v)=x\), then the new observation signature is
\[
o_\Psi(v)=\{q\in Q:q\le x\}=\down x.
\]
The map \(x\mapsto\down x\) is an order embedding of \(Q\) into \(2^Q\), because
\[
x\le y\iff\down x\subseteq\down y.
\]
Surjectivity of \(h\) realizes every such signature and no others. Hence \(P_\Psi\cong Q\).
\end{proof}

\begin{corollary}[Every finite simplicial homotopy type occurs]\label{cor:alltypes}
For every finite simplicial complex \(K\), choose its face poset \(Q=\mathcal X(K)\). Then
\[
\Delta(Q)=\operatorname{sd}(K),
\]
the barycentric subdivision of \(K\). Therefore a refinement above a one-point coarse ProLT can realize the homotopy type of any finite simplicial complex.
\end{corollary}

\begin{theorem}[Undecidability of unrestricted homotopy-neutral refinement]\label{thm:undecidable}
\statusS/\statusA/\statusC\;
There is no algorithm that, for arbitrary finite ProLT observation refinements in unrestricted dimension, always decides whether
\[
\Delta(r_{\Psi\Phi})
\]
is a homotopy equivalence.
\end{theorem}

\begin{proof}
Given a finite simplicial complex \(K\), apply Corollary~\ref{cor:alltypes} to construct a refinement whose coarse signature poset is a point and whose fine order complex has the homotopy type of \(K\). The refinement projection to the point is a homotopy equivalence exactly when \(K\) is contractible. Algorithmic recognition of contractibility for arbitrary finite simplicial complexes is undecidable. Thus a universal decision algorithm for homotopy-neutral ProLT refinements would decide an undecidable problem.
\end{proof}

\begin{corollary}[Homology remains effectively computable]\label{cor:decidablehomology}
For any fixed computable coefficient field, homology-neutrality of a given finite refinement is decidable: construct the finite mapping-cone boundary matrices and compute their ranks. Thus the homotopy and homology classification problems diverge sharply in unrestricted generality.
\end{corollary}

\begin{remark}[Strictness of \(\mathsf H\subsetneq\mathsf M\)]
The universality theorem also imports standard acyclic noncontractible finite complexes. For example, a finite triangulation of a punctured Poincar\'e homology sphere is acyclic but has nontrivial fundamental group. Refining a point to its face poset gives a homology equivalence to the point which is not a homotopy equivalence.
\end{remark}

\section{A practical classification algorithm}

The theory suggests the following release-quality workflow for a concrete observation refinement.

\begin{enumerate}[label=\textbf{Step \arabic*.},leftmargin=4em]
\item \textbf{Build the coarse poset.} Compute the realized old signatures \(P_\Phi\) and their inclusion order.
\item \textbf{Build fiber profiles.} For each \(p\in P_\Phi\), compute
\[
\mathcal B_p=o_\Theta(C_p).
\]
No fine-poset construction is conceptually necessary beyond these profiles.
\item \textbf{Check exact redundancy.} If every \(\mathcal B_p\) is singleton and the profile map is isotone, the refinement is an order isomorphism and the classification stops.
\item \textbf{Check adjoints.} Compute \(\mathsf U(p)\) and \(\mathsf L(p)\). If all \(\mathsf U(p)\in\mathcal B_p\), or all \(\mathsf L(p)\in\mathcal B_p\), the refinement is homotopy-neutral.
\item \textbf{Check principal fibers.} Apply beat-point/collapse tests or another contractibility certificate to \(r^{-1}(\down p)\) and \(r^{-1}(\up p)\). Uniform contractibility certifies simple homotopy equivalence. Uniform acyclicity certifies homology equivalence.
\item \textbf{Compute the exact homology defect.} Build \(\Cone(r_\#)\). Nonzero \(D_k\) proves that the refinement changes homology. Vanishing of all \(D_k\) proves homology equivalence.
\item \textbf{For non-splitting refinements, use descent chains.} Replace the general mapping cone by the relative/descent chain complex of Proposition~\ref{prop:descentcomplex}.
\item \textbf{Do not treat an inconclusive homotopy test as failure.} The Quillen and adjoint criteria are sufficient, not necessary. In unrestricted dimension no complete decision algorithm exists. For a research theorem, restrict to a structurally controlled subclass.
\end{enumerate}

\section{Promising theorem targets after this classification}

The classification above narrows the next research questions substantially.

\begin{enumerate}[label=\textbf{T\arabic*.},leftmargin=2.8em]
\item \textbf{Logical criteria for contractible Quillen fibers.} Find syntactic or semantic conditions on newly added formulas that force every \(r^{-1}(\down p)\) to have a cone point, beat-point reduction, or nerve decomposition.
\item \textbf{Minimal homology-changing observations.} Given \(\Phi\), characterize or optimize observations \(\psi\) for which \(D_k((\Phi,\psi)/\Phi)\neq0\), possibly with a cost on formulas or observations.
\item \textbf{Post-\(T_0\) persistence.} Study the reverse filtration of Corollary~\ref{cor:T0filtration}. Determine whether logical properties of the added observations constrain barcode births/deaths in the relative descent complexes.
\item \textbf{Refinement/update interaction.} Combine the contravariant-looking refinement deletion \(\OX_\Psi\subseteq\OX_\Phi\) after \(T_0\) with premise restriction subcomplexes. This should yield a two-parameter family indexed by observational resolution and admissible-state restriction.
\item \textbf{Cross-layer comparison.} Relate refinement defects in canonical order-complex homology to changes in the presentation-sensitive Dowker compatibility complex and to changes in metric/Rips persistence. A nontrivial comparison theorem would be substantially stronger than defining the layers independently.
\item \textbf{Restricted decidable classes.} Because unrestricted homotopy neutrality is undecidable, identify meaningful ProLT classes where the problem becomes decidable or efficiently certifiable: bounded poset height, bounded number of observations, distributive-lattice images, monotone formula fragments, or bounded treewidth of the comparability/Hasse graph.
\end{enumerate}

\section{Preliminary novelty positioning}

The finite-poset fiber theorems, adjunction-homotopy principle, mapping cones, and undecidability of arbitrary finite-complex contractibility are standard mathematics. They should not be claimed as new. The current contribution candidates are instead the \emph{ProLT-specific synthesis and sharpenings}: the refinement fiber-profile normal form; the exact redundancy theorem in signature terms; the possible/forced monotonicity criterion; the descent-chain description after \(T_0\); the refinement-defect hierarchy; and the universality reduction formulated for propositional observation refinement.

A preliminary literature check found adjacent but distinct areas. Formal concept analysis and information-system theory study binary object--attribute data and attribute reduction. Separately, Gaucher used the phrase ``refinement of observation'' in directed homotopy theory for flows and finite bounded posets, with preservation results for branching/merging homology and underlying homotopy type. That is not the same object as the ProLT signature-projection refinement studied here, but the terminology overlap should be acknowledged in any publication. A dedicated adversarial novelty audit remains necessary before promoting any \statusC\ item to a novelty claim.

\section{Conclusion}

Observation refinement in a ProLT is not a single phenomenon. It may:
\begin{itemize}
\item add a coordinate without changing the topology at all;
\item split observational states while preserving homotopy type;
\item leave all states unsplit but delete specialization comparisons;
\item change homotopy while preserving homology;
\item or change homology itself.
\end{itemize}

The fiber-profile normal form makes these possibilities explicit. It reduces every refinement to a structured surjective poset map
\[
Q_{\mathcal B}\longrightarrow P_\Phi.
\]
Topological redundancy is exactly order isomorphism. Adjoint and Quillen-fiber tests give progressively broader homotopy-neutral classes. Mapping-cone homology gives an exact finite test for homology neutrality. Once the old ProLT is \(T_0\), refinement becomes especially transparent: it can no longer split vertices, so it acts only by deleting chains from the canonical order complex, and the defect is ordinary relative homology generated by the deleted descent chains.

The universality theorem also tells us where to stop looking for an all-purpose homotopy decision rule: unrestricted refinement already contains the full contractibility problem for finite complexes. The productive research direction is therefore not a mythical universal criterion, but a theory of \emph{logically natural subclasses} in which the fiber geometry can be controlled and related back to observations, premises, and eventually time-evolving logical state spaces.

\appendix
\section{Computational verification summary}

The accompanying verifier constructs order complexes and normalized simplicial chain maps explicitly over \(\F\). It verifies:
\begin{enumerate}[label=(\roman*)]
\item the redundant \(X\wedge Y\) example has identical coarse/fine Betti vector \((1,0,0)\) and zero cone homology;
\item the XOR refinement has coarse Betti \((1,0,0)\), fine Betti \((1,0)\), and zero cone homology;
\item the negation split has coarse Betti \((1,0)\), fine Betti \((2)\), and cone Betti \((0,1)\);
\item the splitting-but-neutral \(Y\) refinement of \((X)\) has zero cone homology;
\item the four-chain alternating profile \((1,0,1,0)\) yields contractible coarse and fine complexes while failing both lower and upper Quillen-fiber tests;
\item all \(255\) one-bit profiles over nonempty subposets of \(B_2\) agree with the logical hierarchy reported in Section~\ref{sec:hierarchy}.
\end{enumerate}

\begin{thebibliography}{99}
\bibitem{Theory2026}
B. Theory,
\emph{Propositional Logical Topology: Observation, Distinguishability, and Inference},
core manuscript v0.8, September 2026, unpublished.

\bibitem{TheoryHomology2026}
B. Theory,
\emph{Homological Foundations for Propositional Logical Topologies},
working research note v0.1, September 2026, unpublished.

\bibitem{McCord1966}
M. C. McCord,
``Singular homology groups and homotopy groups of finite topological spaces,''
\emph{Duke Mathematical Journal} 33 (1966), 465--474.

\bibitem{Quillen1978}
D. Quillen,
``Homotopy properties of the poset of nontrivial \(p\)-subgroups of a group,''
\emph{Advances in Mathematics} 28 (1978), 101--128.

\bibitem{Barmak2011}
J. A. Barmak,
``On Quillen's Theorem A for posets,''
\emph{Journal of Combinatorial Theory, Series A} 118 (2011), 2445--2453.

\bibitem{BarmakBook2011}
J. A. Barmak,
\emph{Algebraic Topology of Finite Topological Spaces and Applications},
Lecture Notes in Mathematics 2032, Springer, 2011.

\bibitem{Tancer2016}
M. Tancer,
``Recognition of collapsible complexes is NP-complete,''
\emph{Discrete \& Computational Geometry} 55 (2016), 21--38.
The appendix discusses undecidability of contractibility for finite simplicial complexes.

\bibitem{Gaucher2006}
P. Gaucher,
``T-homotopy and refinement of observation. III. Invariance of the branching and merging homologies,''
\emph{New York Journal of Mathematics} 12 (2006), 319--348.

\bibitem{Gaucher2005}
P. Gaucher,
``T-homotopy and refinement of observation. IV. Invariance of the underlying homotopy type,''
arXiv:math/0505331 (2005).

\end{thebibliography}

\end{document}
